\subsection{1.1 Source basis and domains}\label{source-basis-and-domains}

The two newly supplied files are preserved unchanged under \texttt{inputs/}. The mathematical input used here is their one-factor coercivity P1--P11, joint minimum J1--J8, spectral comparison S1--S5, and, on the simple-quartet domain, V1--V6 and V13. The EIQ monic-profile theorem in V1 is an imported written source theorem; the present note does not independently prove that asymptotic. The relevant GitHub source was read at commit \nolinkurl{ff251e3adc5fd78049170b8fc94c26f882df3d14}, file \nolinkurl{workbenches/splitzero-tandem/continuations/20260920-original-spectral-determinant/ORIGINAL_SPECTRAL_DETERMINANT.tex}, blob \nolinkurl{11a2154819682c01c878ffb0b886f7b704fed888}. Its SD1--SD12 specify the original companion, the selfadjoint source compression, the rank-one correction, and its singular-value bounds. SD26--SD28 keep the original metric dependence. We use those original definitions; finite identities needed below are also proved here.

The evaluated heat threshold is on the programme's stated simple quartet: fixed h, fixed delta in (0,1/2), gamma\textgreater2, k=1 mod 4, and q=(k+1)\^{}2. This is not a claim to have exhibited an off-critical xi zero. The general covariance identities in the later parts require neither a quartet nor equal multiplicities.

Let chi be the whole original sum polynomial and let M be multiplication by y in E=C{[}y{]}/chi\_y, with y=(S-k/2)/i and chi\_y(y)=i\^{}(-q)chi(k/2+iy). Let G\_N be the original attained metric from the actual convolution m\_k=w\_h\^{}\{*k\}. Write C\_N for its selfadjoint physical y-compression. Then \[
 M=C_N+R_N,\quad R_N=r_N\ell_N,\quad R_N^2=0,\quad
 C_N^{\dagger_{G_N}}=C_N,\quad \epsilon_N=\|R_N\|_{G_N}>0.
 \tag{H1}
\] The sign is the SD2 convention. It agrees with an earlier M=A-Q convention by R=-Q; the positive scalar epsilon is unchanged.

Put r=N+1-q and s\_N=r/q. We use q-1\textless=N\textless=2q, including the fourth original determinant endpoint. The finite multiplication estimate applies at degree 2q as well as at 2q-1, so one fixed constant B\_h gives \[
 \|C_N\|_{G_N}\le B_hq.
 \tag{H2}
\] When a source moment above degree 2q is used, the supplied all-degree coercivity is applied through degree 2q+2. This retains the degree-raising column; it does not silently use a truncated multiplication form.

\subsection{1.2 The entire allowance profile follows from the retained monic minima}\label{the-entire-allowance-profile-follows-from-the-retained-monic-minima}

Let omega\_n be the squared norm of the original real monic source polynomial p\_n.~Let nu\_j be the squared monic norm for the complete relation measure chi\_y\^{}2 m\_k. Define \[
 d_n=\omega_n/\nu_{n-q},\qquad b_{n+1}=\omega_{n+1}/\omega_n.
\] Direct projection onto the last monic relation proves \[
 \epsilon_{q-1}^2=b_q(1-d_q)/d_q,
\] \[
 \epsilon_N^2=b_{N+1}(1-d_N)(1-d_{N+1})/d_{N+1}\quad(N\ge q).
 \tag{H3}
\] Here is the full norm calculation. At cutoff N, the minimum representative of {[}p\_(N+1){]} is the difference between the monic relation of degree N+1 and p\_(N+1), up to sign. Its squared norm is nu\_r-omega\_(N+1). The leading-coefficient functional on the minimum source has squared dual norm (1-d\_N)/omega\_N for N\textgreater=q, and 1/omega\_(q-1) at the first cutoff. Indeed {[}p\_N{]} has squared quotient norm omega\_N-omega\_N\^{}2/nu\_(N-q), and the leading functional is its inner product divided by omega\_N. Multiplying these two norm factors gives H3. Every preceding relation is present in these orthogonal projections.

The supplied V4 gives, through r=q+1, \[
 \log(\nu_r/\omega_{q+r})=2q\psi(r/q)+O_h(k\log^2(q+2)).
 \tag{H4}
\] The same full-source comparisons applied to two consecutive monic affine fibres give log b\_(N+1)=O\_h(k log\^{}2(q+2)). The source profile is continuous and decreasing on the required interval, with psi(1)\textgreater0. Consequently d\_n is uniformly exponentially small in q on this window; both logarithms log(1-d\_n) in H3 are retained and tend to zero. We obtain \[
 \boxed{\log\epsilon_N^2=2q\psi(s_N)+O_h(k\log^2(q+2)).}
 \tag{H5}
\] This is uniform in N. It is an asymptotic with the original profile theorem's eventual domain, not an effective numerical threshold in k.

\subsection{1.3 An exact elliptic expression for the two heat thresholds}\label{an-exact-elliptic-expression-for-the-two-heat-thresholds}

Denote the complete elliptic integrals by K(u), E(u), with modulus u in (0,1), and put u'=sqrt(1-u\^{}2). The supplied equilibrium endpoint equations and profile derivative are \[
 1+s=E(u)/(u'K(u)),\qquad \psi'(s)=\log(u/E(u)),\qquad \psi(0)=\log(4/\pi).
 \tag{H6}
\] The parameter map in H6 is invertible: its derivative is \[
 \frac{(K-E)(E-u'^2K)}{u\,u'^3K^2}>0.
\] Both factors in the numerator are positive by the defining elliptic integrals. Its range is (0,infinity) after subtracting one. The two derivative formulas E'=(E-K)/u and K'=E/(u u'\^{}2)-K/u follow by differentiation of their defining integrals.

Integration by parts in H6 gives the complete closed expression \[
 \boxed{\psi(s)=s\log\frac{u_s}{E(u_s)}+
                   \log\frac{1+u_s'}{E(u_s)},\quad
       1+s=\frac{E(u_s)}{u_s'K(u_s)}.}
 \tag{H7}
\] For a direct verification, differentiate the right side: the coefficient of du/ds is ((1+s)K/E-1)/u-u/(u'(1+u'))=0. Its limit at zero is log(4/pi). Thus this is the same profile, through its exact inverse parameter map.

Let a\_0=psi(0), a\_1=psi(1). Then \[
 a_0=\log(4/\pi),\qquad
 a_1=\log\frac{u_1(1+u_1')}{E(u_1)^2},\qquad
 E(u_1)=2u_1'K(u_1).
 \tag{H8}
\] The exact rational enclosure in \nolinkurl{checks/heat_checks.py} proves \[
 \boxed{0<a_1<1/8<a_0.}\tag{H9}
\] For completeness, the enclosure uses u\_1 in (987/1000,988/1000). Write z=u\^{}2 and c\_n=(binom(2n,n)/4\textsuperscript{n)}2. The series K=(pi/2)sum c\_n z\^{}n and E=(pi/2)(1-sum\_(n\textgreater=1)c\_n z\^{}n/(2n-1)) have omitted positive tails bounded by c\_(J+1)z\^{}(J+1)/(1-z) and this quantity/(2J+1). At J=512 the rational squared inequalities E\_series\^{}2 versus 4(1-z)K\_series\^{}2 prove the root bracket. Using 333/106\textless pi\textless355/113 and 154/1000\textless u\_1'\textless161/1000 then proves 1\textless u\_1(1+u\_1')/E(u\_1)\^{}2\textless9/8. Finally log(9/8)\textless1/8 and exp(1/8)\textless=8/7\textless4/(355/113). The checker separately derives the stated loose pi bounds from Machin's alternating arctangent series.

Non-interval high-precision evaluations, not native moment evaluations, are \[
 u_1\simeq0.9871005293061829185,\quad
 a_1\simeq0.06651895202027391967,\quad
 a_0\simeq0.24156447527049044469.
\]
